%========================
% chapters/09_conclusion.tex
%========================
\chapter{Conclusion}
\label{chap:conclusion}

This dissertation studied node ordering as a central component of autoregressive graph generation. The starting point was a simple observation: sequential graph generators do not observe unordered graphs directly. They observe construction sequences, and these sequences depend on how the nodes are ordered. Therefore, an ordering is not just a preprocessing convention; it defines the conditional histories, action candidates, and structural regularities available to the generator during training.

The first contribution was a controlled analysis of deterministic orderings under a DGMG-style generator. By keeping the generator fixed and changing only the graph-to-sequence map, the dissertation isolated the effect of ordering. The results showed that orderings induce strong and topology-dependent differences. Cycles favor path-like orderings such as DFS and LexBFS. Trees and binary trees favor BFS because it preserves parent-frontier locality. Stars favor degree ordering because early hub exposure simplifies most destination decisions. Barab\'asi--Albert graphs reveal a more complex trade-off, in which different orderings emphasize size, connectivity, degree inequality, or hub dominance. These findings support the interpretation that ordering acts as an entropy-shaping mechanism for the sequential decision problem.

The second and main contribution was a learned-ordering framework based on reinforcement learning. The dissertation formulated node sequencing as a graph-conditioned Markov Decision Process in which a policy selects one unselected node at a time until a complete ordering is produced. The resulting ordering is translated into a DGMG action sequence, and the reward is computed from downstream generator behavior. The reward design explicitly combines predictive fit, measured as negative validation log-likelihood, with structural validity metrics computed on freshly sampled generated graphs. Because the ordering decisions are discrete and their effect is delayed, the policy is optimized with a policy-gradient estimator. The theorem presented in Chapter~\ref{chap:rl-ordering} formalizes this gradient and justifies the update used by the method.

An intermediate GCN-based RL policy confirmed that the general learned-ordering framework is viable and capable of achieving strong size control on BA graphs. However, connectivity remained a weakness, motivating the architectural upgrades introduced in the final configuration. The final version of the method, Attention RL ordering, strengthened the policy with an attention-based graph encoder, three graph layers, and a degree-based warm start. The method was evaluated on Barab\'asi--Albert graphs, where it achieved perfect valid-size ratio, near-perfect connectivity, realistic average size, and competitive hubiness. The result is significant because it combines properties that are separated across deterministic baselines: BFS is strong for connectivity but weak for size, while degree ordering is strong for size and hub exposure but weak for connectivity. The learned method moves toward a hybrid region of the metric space, which is precisely the motivation for replacing a fixed heuristic with an adaptive policy.

The experiments are based on synthetic graph families, and although these families are useful for controlled interpretation, real-world graph distributions may introduce additional constraints such as labels, attributes, edge types, and domain-specific validity. The learned-ordering method also has higher computational cost than deterministic ordering and depends on reward design. Policy-gradient optimization can be noisy, especially when the reward is based on generator training and sample evaluation. These limitations do not undermine the main conclusion, but they define the scope in which the results should be generalized.

Several directions remain open. The first is to evaluate learned ordering on real datasets such as molecular graphs, program graphs, or biological networks. The second is to combine learned ordering with stronger autoregressive backbones, including transformer-based graph generators such as GraphGPT \cite{zhao2023graphgpt} and G2PT \cite{chen2025g2pt}. The third is to improve the reinforcement-learning component with variance reduction, actor-critic methods, imitation from strong heuristics, or curriculum learning over graph size. The fourth is to study whether the learned policy discovers interpretable hybrid strategies, for example by measuring agreement with BFS, degree ordering, or MCS at different stages of the trajectory. The policy-entropy diagnostic shown in Figure~\ref{fig:rl-entropy-mean} suggests that the policy becomes increasingly concentrated over training, which is consistent with a transition from broad exploration to confident ordering decisions; characterizing when and why this concentration occurs could yield useful insights for improving training stability.

Overall, the dissertation argues that the ordering problem should be treated as part of the model design in sequential graph generation. A graph generator trained on ordered sequences is only as effective as the representation through which the graph is exposed. Handcrafted orderings provide useful priors, but they are topology-dependent and metric-dependent. Learned orderings offer a way to adapt this prior to the data and to the downstream generator. This perspective opens a path toward graph generators in which the construction order is not assumed in advance, but learned as part of the generative process itself.
